% ===========================================================================
%  Chapitre 9 — Arithmétique dans Z
%  PE 2026, composante ALGÈBRE
% ===========================================================================
\bmtchapitre{Arithmétique dans l'ensemble des entiers}
  {Résoudre des exercices et des problèmes d'arithmétique : divisibilité,
   division euclidienne, congruences, PGCD et PPCM, théorèmes de Bézout et de
   Gauss, équations diophantiennes, ensemble des classes de congruence.}
  {Algèbre \textbullet\ environ 14 heures}

L'arithmétique porte sur les entiers, c'est-à-dire sur les objets
mathématiques les plus familiers qui soient. C'est précisément ce qui la rend
déroutante : les énoncés se comprennent en une seconde, les démonstrations
demandent une heure. On ne calcule presque pas dans ce chapitre, on raisonne.

Vous verrez que l'outil décisif tient en une phrase : dans une division, le
reste dit tout. Deux entiers qui laissent le même reste se comportent de la
même façon vis-à-vis du diviseur, et cette remarque suffit à établir des
divisibilités qu'aucun calcul direct n'atteindrait. C'est l'idée des
congruences, et le baccalauréat les place presque chaque année en ouverture de
sujet.

% ---------------------------------------------------------------------------
\begin{revision}
Tout ce qui suit a été vu au collège ou en classe de première, sauf la
dernière question, qui reprend le principe de récurrence du chapitre précédent.

\begin{enumerate}[label=\textbf{\arabic*.}, leftmargin=*, itemsep=0.35em]
  \item Effectuer la division euclidienne de $47$ par $6$, puis de $-47$ par
        $6$.
  \item Décomposer $360$ et $84$ en produits de facteurs premiers.
  \item Donner la liste des diviseurs positifs de $36$.
  \item Rappeler les critères de divisibilité par $2$, $3$, $5$, $9$ et $11$.
  \item Qu'appelle-t-on un nombre premier ? Les premiers nombres premiers.
  \item Démontrer par récurrence que $\displaystyle\sum_{k=1}^{n}k
        = \dfrac{n(n+1)}{2}$.
\end{enumerate}
\end{revision}

\begin{automatismes}
Sans calculatrice, sans brouillon.

\begin{enumerate}[label=\textbf{\alph*)}, leftmargin=*, itemsep=0.25em]
  \item Reste de $58$ dans la division par $7$
  \item Reste de $100$ dans la division par $9$
  \item Diviseurs positifs de $18$
  \item $360 = 2^{?}\times3^{?}\times5^{?}$
  \item $1001$ est-il divisible par $7$ ?
  \item Le plus grand diviseur commun de $24$ et $36$
  \item Le plus petit multiple commun de $6$ et $8$
  \item $2^{10}$
\end{enumerate}
\end{automatismes}

\begin{corrige}[automatismes]
a) $2$ \quad b) $1$ \quad c) $1,2,3,6,9,18$ \quad
d) $2^{3}\times3^{2}\times5$ \quad e) oui, $1001 = 7\times143$ \quad
f) $12$ \quad g) $24$ \quad h) $1024$.
\end{corrige}

% ===========================================================================
\section{Divisibilité dans l'ensemble des entiers relatifs}

\begin{definition}[diviseur, multiple]
Soient $a$ et $b$ deux entiers relatifs. On dit que $b$ \textbf{divise} $a$,
et l'on écrit $b \divi a$, s'il existe un entier relatif $k$ tel que
$a = kb$. On dit alors que $a$ est un \textbf{multiple} de $b$, et que $b$ est
un \textbf{diviseur} de $a$.
\end{definition}

\begin{remarque}
Tout entier divise $0$, car $0 = 0 \times b$. En revanche $0$ ne divise que
lui-même. Et $1$ divise tout entier. Ces trois remarques sont des cas
particuliers que l'on oublie systématiquement en début d'année ; il vaut mieux
les fixer tout de suite.
\end{remarque}

\begin{propriete}[règles de calcul]
Soient $a$, $b$, $c$ des entiers relatifs.
\begin{enumerate}[label=\textbf{\arabic*.}, leftmargin=*, itemsep=0.15em]
  \item Si $b \divi a$ et $a \divi c$, alors $b \divi c$.
  \item Si $b \divi a$ et $b \divi c$, alors $b$ divise $au+cv$ pour tous
        entiers $u$ et $v$.
  \item Si $b \divi a$ et $a \neq 0$, alors $\abs{b} \leq \abs{a}$.
\end{enumerate}
\end{propriete}

\begin{demonstration}
\textbf{1.} $a = kb$ et $c = k'a$ donnent $c = k'kb$.

\textbf{2.} $a = kb$ et $c = k'b$ donnent $au+cv = (ku+k'v)\,b$.

\textbf{3.} $a = kb$ avec $a \neq 0$ impose $k \neq 0$, donc
$\abs{k} \geq 1$ et $\abs{a} = \abs{k}\abs{b} \geq \abs{b}$.
\end{demonstration}

\begin{avertissement}
La troisième propriété est la seule qui permette de \emph{conclure} qu'un
ensemble de diviseurs est fini. C'est elle qu'on invoque quand un exercice
demande de trouver tous les entiers $n$ tels que $n+2$ divise $3n+7$ : on
fabrique une combinaison qui élimine $n$, et le diviseur se retrouve
prisonnier d'un nombre fixe.
\end{avertissement}

\begin{exemple}[éliminer l'inconnue]
Cherchons les entiers relatifs $n$ tels que $n-3$ divise $n+5$.

Puisque $n-3$ divise $n-3$, il divise aussi la différence
$(n+5)-(n-3) = 8$. Donc $n-3$ appartient à l'ensemble des diviseurs de $8$,
soit $\{-8,-4,-2,-1,1,2,4,8\}$, ce qui donne
\[
  n \in \{-5,\ -1,\ 1,\ 2,\ 4,\ 5,\ 7,\ 11\}.
\]
Réciproquement, chacune de ces valeurs convient : la condition trouvée est
équivalente à celle de départ, puisque nous n'avons utilisé que des
équivalences.
\end{exemple}

\begin{methode}{montrer qu'un entier variable en divise un autre}
Trois gestes, toujours les mêmes.
\begin{enumerate}[label=\textbf{\arabic*.}, leftmargin=*, itemsep=0.15em]
  \item Écrire ce que l'on sait sous la forme « $d$ divise \dots ».
  \item Combiner ces divisibilités pour \emph{éliminer} la lettre variable.
  \item Conclure que $d$ divise un entier fixe, puis lister ses diviseurs.
\end{enumerate}
\end{methode}

\begin{entrainement}[divisibilité]
\exo[\niveau{1}] Déterminer tous les diviseurs positifs de $84$.

\exo[\niveau{1}] Montrer que si $n$ est un entier, $n(n+1)$ est pair.

\exo[\niveau{2}] Trouver tous les entiers relatifs $n$ tels que $n+1$ divise
$n^{2}+3$.

\exo[\niveau{2}] Montrer que pour tout entier $n$, $6$ divise
$n^{3}-n$.

\exo[\niveau{3}] Soient $a$ et $b$ deux entiers. Montrer que si $a \divi b$ et
$b \divi a$, alors $a = b$ ou $a = -b$.
\end{entrainement}

\begin{corrige}
\textbf{1.} $84 = 2^{2}\times3\times7$, d'où les douze diviseurs positifs
$1$, $2$, $3$, $4$, $6$, $7$, $12$, $14$, $21$, $28$, $42$, $84$.

\textbf{2.} Parmi deux entiers consécutifs, l'un est pair : leur produit
l'est donc aussi.

\textbf{3.} $n^{2}+3 = (n+1)(n-1)+4$. Comme $n+1$ divise le premier terme,
il divise $n^{2}+3$ si et seulement s'il divise $4$, c'est-à-dire si
$n+1 \in \{-4\,;-2\,;-1\,;1\,;2\,;4\}$, soit
$n \in \{-5\,;-3\,;-2\,;0\,;1\,;3\}$.

\textbf{4.} $n^{3}-n = (n-1)n(n+1)$ est le produit de trois entiers
consécutifs : l'un au moins est pair, l'un exactement est multiple de $3$.
Le produit est donc divisible par $2$ et par $3$, donc par $6$.

\textbf{5.} Si $a = 0$, alors $a \divi b$ impose $b = 0$ et la conclusion
vaut. Sinon, $a \divi b$ donne $\abs a \leq \abs b$ et $b \divi a$ donne
$\abs b \leq \abs a$ : les deux valeurs absolues sont égales, donc
$a = b$ ou $a = -b$.
\end{corrige}

% ===========================================================================
\section{Division euclidienne des entiers}

\begin{theoreme}[division euclidienne]
Soient $a$ un entier relatif et $b$ un entier naturel non nul. Il existe un
unique couple $(q\,;r)$ d'entiers, avec $q \in \Z$ et $r \in \N$, tel que
\[
  a = bq+r \qquad \text{et} \qquad 0 \leq r < b .
\]
L'entier $q$ est le \textbf{quotient}, l'entier $r$ le \textbf{reste}.
\end{theoreme}

\begin{demonstration}
\emph{Existence.} L'ensemble des multiples de $b$ inférieurs ou égaux à $a$
est non vide et majoré par $a$ ; il possède donc un plus grand élément $bq$.
On pose $r = a-bq \geq 0$. Si l'on avait $r \geq b$, alors $b(q+1) \leq a$,
ce qui contredirait le choix de $q$. Donc $0 \leq r < b$.

\smallskip
\emph{Unicité.} Si $a = bq+r = bq'+r'$ avec $0 \leq r,r' < b$, alors
$b(q-q') = r'-r$. Or $\abs{r'-r} < b$, donc $b$ divise un entier de valeur
absolue strictement inférieure à $b$ : cet entier est nul. D'où $r = r'$ puis
$q = q'$.
\end{demonstration}

\begin{visualisation}[le reste, lu sur une règle graduée]
\begin{center}
\begin{tikzpicture}[x=1.05cm, y=1cm]
  \draw[bmtGris, line width=0.8pt, -{Stealth[length=3mm]}] (-0.4,0) -- (9.6,0);
  \foreach \k in {0,1,2,3,4,5,6,7,8,9}
    \draw[bmtGris] (\k,0.1) -- (\k,-0.1);
  \foreach \k/\lab in {0/0, 7/7} {
    \draw[bmtLaterite, line width=1.5pt] (\k,0.16) -- (\k,-0.16);
    \node[bmtLaterite, font=\footnotesize, anchor=north] at (\k,-0.22) {$\lab$};}
  \node[bmtLaterite, font=\footnotesize, anchor=north] at (3.5,-0.22) {};
  \draw[bmtIndigo, line width=1.4pt] (7,0.55) -- (9,0.55);
  \node[bmtIndigo, font=\footnotesize, anchor=south] at (8,0.6) {reste $= 2$};
  \draw[bmtVetiver, line width=1.4pt] (0,0.55) -- (7,0.55);
  \node[bmtVetiver, font=\footnotesize, anchor=south] at (3.5,0.6)
    {$q$ fois $b$, ici $1 \times 7$};
  \filldraw[bmtEncre] (9,0) circle (2pt);
  \node[bmtEncre, font=\footnotesize, anchor=north] at (9,-0.22) {$a = 9$};
\end{tikzpicture}
\end{center}

Poser la division euclidienne de $9$ par $7$, c'est avancer sur la droite par
bonds de longueur $7$ aussi longtemps qu'on ne dépasse pas $9$, puis regarder
ce qu'il reste à parcourir. Le reste est toujours \emph{en avant} du dernier
bond, et strictement plus court qu'un bond : voilà pourquoi
$0 \leq r < b$. Pour $a = -9$, le raisonnement est le même en partant de la
gauche : les bonds vont jusqu'à $-14$, et le reste vaut $5$.
\end{visualisation}

\begin{avertissement}
Le reste est toujours \emph{positif ou nul}, même lorsque $a$ est négatif. La
division de $-9$ par $7$ s'écrit $-9 = 7\times(-2)+5$, et non
$-9 = 7\times(-1)-2$. Une calculatrice qui affiche $-2$ répond à une autre
question que la nôtre.
\end{avertissement}

\begin{entrainement}[division euclidienne]
\exo[\niveau{1}] Effectuer la division euclidienne de $2026$ par $17$.

\exo[\niveau{1}] Effectuer la division euclidienne de $-100$ par $9$, puis de
$-100$ par $7$.

\exo[\niveau{2}] Le reste de la division de $a$ par $12$ est $7$. Quel est le
reste de la division de $a$ par $4$ ? Et par $3$ ?

\exo[\niveau{3}] Déterminer les entiers naturels $a$ dont la division
euclidienne par $8$ donne un quotient égal au reste.
\end{entrainement}

\begin{corrige}
\textbf{6.} $2026 = 17\times119+3$, avec $0 \leq 3 < 17$.

\textbf{7.} $-100 = 9\times(-12)+8$ et $-100 = 7\times(-15)+5$. Le reste est
toujours positif : c'est le quotient qui absorbe le signe.

\textbf{8.} $a = 12q+7$. Comme $12q$ est multiple de $4$ et que $7 = 4+3$, le
reste modulo $4$ vaut $3$. Comme $12q$ est multiple de $3$ et que
$7 = 2\times3+1$, le reste modulo $3$ vaut $1$.

\textbf{9.} La condition s'écrit $a = 8q+q = 9q$ avec $0 \leq q < 8$ : les
entiers cherchés sont $0$, $9$, $18$, $27$, $36$, $45$, $54$ et $63$.
\end{corrige}

% ===========================================================================
\section{Congruences et leurs propriétés}

\begin{definition}[congruence modulo n]
Soit $n$ un entier naturel non nul. Deux entiers relatifs $a$ et $b$ sont
\textbf{congrus modulo} $n$ lorsque $n$ divise $a-b$. On écrit
\[
  \congru{a}{b}{n} .
\]
\end{definition}

\begin{propriete}[caractérisation par le reste]
$\congru{a}{b}{n}$ si et seulement si $a$ et $b$ ont le même reste dans la
division euclidienne par $n$.
\end{propriete}

\begin{demonstration}
Écrivons $a = nq+r$ et $b = nq'+r'$ avec $0 \leq r,r' < n$. Alors
$a-b = n(q-q')+(r-r')$, donc $n \divi a-b$ équivaut à $n \divi r-r'$. Or
$\abs{r-r'} < n$ : cette divisibilité équivaut à $r = r'$.
\end{demonstration}

\begin{propriete}[compatibilité avec les opérations]
Si $\congru{a}{b}{n}$ et $\congru{c}{d}{n}$, alors
\[
  \congru{a+c}{b+d}{n}, \qquad \congru{ac}{bd}{n}, \qquad
  \congru{a^{k}}{b^{k}}{n} \ \text{ pour tout } k \in \N .
\]
\end{propriete}

\begin{demonstration}
Pour la somme, $(a+c)-(b+d) = (a-b)+(c-d)$ est somme de deux multiples de
$n$. Pour le produit, on écrit
\[
  ac-bd = ac-bc+bc-bd = c(a-b)+b(c-d),
\]
somme de deux multiples de $n$. La puissance s'en déduit par récurrence sur
$k$ : vraie pour $k = 0$, et si $\congru{a^{k}}{b^{k}}{n}$, le produit par
$\congru{a}{b}{n}$ donne $\congru{a^{k+1}}{b^{k+1}}{n}$.
\end{demonstration}

\begin{avertissement}
On ne simplifie pas une congruence comme une égalité. De
$\congru{2\times3}{2\times0}{6}$ on ne tire pas $\congru{3}{0}{6}$. La
simplification par $c$ n'est licite que si $c$ et $n$ sont premiers entre eux
— c'est une conséquence du théorème de Gauss, établi plus loin.
\end{avertissement}

\begin{exemple}[une divisibilité qu'aucun calcul direct n'atteint]
Montrons que pour tout entier naturel $n$, l'entier
$N = 9^{3n+1}+4^{6n+1}$ est divisible par $13$.

On travaille modulo $13$. D'abord
$9^{3} = 729 = 13\times56+1$, donc $\congru{9^{3}}{1}{13}$, d'où
$\congru{9^{3n}}{1}{13}$ et $\congru{9^{3n+1}}{9}{13}$.

Ensuite $4^{6} = 4096 = 13\times315+1$, donc $\congru{4^{6}}{1}{13}$, puis
$\congru{4^{6n+1}}{4}{13}$.

En additionnant, $\congru{N}{9+4}{13}$, c'est-à-dire
$\congru{N}{13}{13}$, donc $\congru{N}{0}{13}$ : l'entier $N$ est divisible
par $13$.
\end{exemple}

\begin{methode}{établir une divisibilité dépendant de n}
On cherche d'abord la plus petite puissance de la base qui vaut $1$ modulo
$n$. Toute puissance multiple de celle-là vaut alors $1$, et il ne reste qu'un
facteur constant à traiter. C'est la méthode attendue au baccalauréat, et elle
remplace avantageusement une récurrence — que l'on peut néanmoins produire si
l'énoncé l'exige.
\end{methode}

\begin{entrainement}[congruences]
\exo[\niveau{1}] Déterminer le reste de $2^{100}$ dans la division par $7$.

\exo[\niveau{1}] Résoudre dans $\Z$ la congruence $\congru{3x}{1}{5}$.

\exo[\niveau{2}] Montrer que pour tout entier $n$, $\congru{n^{2}}{0}{4}$ ou
$\congru{n^{2}}{1}{4}$. En déduire qu'aucun carré ne vaut $3$ modulo $4$.

\exo[\niveau{2}] Montrer que $4^{n}+6n-1$ est divisible par $9$ pour tout
entier naturel $n$.

\exo[\niveau{3}] Quel est le chiffre des unités de $7^{2026}$ ?
\end{entrainement}

\begin{corrige}
\textbf{10.} $2^{3} = 8 \equiv 1 \pmod 7$, et $100 = 3\times33+1$, donc
$2^{100} = \left(2^{3}\right)^{33}\times2 \equiv 2 \pmod 7$ : le reste vaut
$2$.

\textbf{11.} On essaie les cinq classes : $3\times2 = 6 \equiv 1$. Les
solutions sont les entiers $x$ tels que $\congru{x}{2}{5}$.

\textbf{12.} Si $n = 2k$, $n^{2} = 4k^{2} \equiv 0$ ; si $n = 2k+1$,
$n^{2} = 4\left(k^{2}+k\right)+1 \equiv 1$. Un carré est donc congru à $0$ ou
à $1$ modulo $4$, jamais à $3$.

\textbf{13.} Au rang $0$ : $1+0-1 = 0$, divisible par $9$. Si
$4^{n}+6n-1 = 9m$, alors
\[
  4^{n+1}+6(n+1)-1 = 4\left(4^{n}+6n-1\right)-18n+9 = 9(4m-2n+1),
\]
qui est encore divisible par $9$.

\textbf{14.} Les chiffres des unités de $7^{n}$ se répètent selon le cycle
$7$, $9$, $3$, $1$, de longueur $4$. Comme $2026 = 4\times506+2$, le chiffre
cherché est le deuxième du cycle, soit $9$.
\end{corrige}

% ===========================================================================
\section{Le plus grand diviseur commun et le plus petit multiple commun}

\begin{definition}[PGCD, PPCM]
Soient $a$ et $b$ deux entiers non tous deux nuls. Le \textbf{PGCD} de $a$ et
$b$, noté $\pgcd(a\,;b)$, est le plus grand entier naturel qui divise à la
fois $a$ et $b$. Le \textbf{PPCM}, noté $\ppcm(a\,;b)$, est le plus petit
entier naturel non nul multiple à la fois de $a$ et de $b$.
\end{definition}

\begin{theoreme}[algorithme d'Euclide]
Si $a = bq+r$, alors $\pgcd(a\,;b) = \pgcd(b\,;r)$. En répétant la division
jusqu'à obtenir un reste nul, le dernier reste non nul est le PGCD.
\end{theoreme}

\begin{demonstration}
Tout diviseur commun à $a$ et $b$ divise $r = a-bq$ : c'est donc un diviseur
commun à $b$ et $r$. Réciproquement, tout diviseur commun à $b$ et $r$ divise
$a = bq+r$. Les deux couples ont exactement les mêmes diviseurs communs, donc
le même plus grand. Les restes successifs forment une suite d'entiers naturels
strictement décroissante : elle atteint $0$ en un nombre fini d'étapes.
\end{demonstration}

\begin{exemple}[l'algorithme déroulé]
Calculons $\pgcd(1071\,;462)$.
\[
\begin{aligned}
  1071 &= 2\times462+147, \\
  462  &= 3\times147+21, \\
  147  &= 7\times21+0 .
\end{aligned}
\]
Le dernier reste non nul est $21$ : $\pgcd(1071\,;462) = 21$.

\smallskip
En remontant les égalités, on obtient de surcroît une relation de Bézout :
$21 = 462-3\times147 = 462-3(1071-2\times462) = 7\times462-3\times1071$.
\end{exemple}

\begin{propriete}[relation entre PGCD et PPCM]
Pour tous entiers naturels non nuls $a$ et $b$,
\[
  \pgcd(a\,;b)\times\ppcm(a\,;b) = ab .
\]
\end{propriete}

\begin{entrainement}[PGCD et PPCM]
\exo[\niveau{1}] Calculer $\pgcd(360\,;84)$ par l'algorithme d'Euclide, puis
$\ppcm(360\,;84)$.

\exo[\niveau{2}] Déterminer les couples d'entiers naturels $(a\,;b)$ tels que
$\pgcd(a\,;b) = 12$ et $a+b = 96$.

\exo[\niveau{2}] Écrire une relation de Bézout entre $17$ et $24$.

\exo[\niveau{3}] Déterminer les couples $(a\,;b)$ d'entiers naturels vérifiant
$\pgcd(a\,;b) = 6$ et $\ppcm(a\,;b) = 180$.
\end{entrainement}

\begin{corrige}
\textbf{15.} $360 = 4\times84+24$, $84 = 3\times24+12$, $24 = 2\times12$ :
le dernier reste non nul est $12$, donc $\pgcd(360\,;84) = 12$. Puis
$\ppcm(360\,;84) = \dfrac{360\times84}{12} = 2520$.

\textbf{16.} On pose $a = 12a'$ et $b = 12b'$ avec $\pgcd(a'\,;b') = 1$. La
condition $a+b = 96$ devient $a'+b' = 8$, et les couples d'entiers premiers
entre eux de somme $8$ sont $(1\,;7)$, $(3\,;5)$, $(5\,;3)$ et $(7\,;1)$.
D'où $(a\,;b) \in \{(12\,;84), (36\,;60), (60\,;36), (84\,;12)\}$.

\textbf{17.} L'algorithme d'Euclide donne $24 = 17+7$, $17 = 2\times7+3$,
$7 = 2\times3+1$. En remontant,
$1 = 7-2\times3 = 5\times7-2\times17 = 5\times24-7\times17$, soit
$17\times(-7)+24\times5 = 1$.

\textbf{18.} On pose $a = 6a'$ et $b = 6b'$ avec $\pgcd(a'\,;b') = 1$. Alors
$\ppcm(a\,;b) = 6a'b' = 180$, donc $a'b' = 30$. Les couples d'entiers premiers
entre eux de produit $30$ sont $(1\,;30)$, $(2\,;15)$, $(3\,;10)$, $(5\,;6)$
et leurs symétriques. D'où
$(a\,;b) \in \{(6\,;180), (12\,;90), (18\,;60), (30\,;36)\}$, et les couples
obtenus en échangeant $a$ et $b$.
\end{corrige}

% ===========================================================================
\section{Nombres premiers et nombres premiers entre eux}

\begin{definition}[nombre premier, entiers premiers entre eux]
Un entier naturel $p \geq 2$ est \textbf{premier} lorsque ses seuls diviseurs
positifs sont $1$ et $p$. Deux entiers $a$ et $b$ sont \textbf{premiers entre
eux} lorsque $\pgcd(a\,;b) = 1$.
\end{definition}

\begin{theoreme}[décomposition en facteurs premiers]
Tout entier naturel $n \geq 2$ s'écrit de manière unique, à l'ordre des
facteurs près, sous la forme
$n = p_{1}^{\alpha_{1}}p_{2}^{\alpha_{2}}\cdots p_{k}^{\alpha_{k}}$, où les
$p_{i}$ sont des nombres premiers deux à deux distincts et les $\alpha_{i}$
des entiers naturels non nuls.
\end{theoreme}

\begin{propriete}[nombre de diviseurs]
Si $n = p_{1}^{\alpha_{1}}\cdots p_{k}^{\alpha_{k}}$, le nombre de diviseurs
positifs de $n$ vaut $(\alpha_{1}+1)(\alpha_{2}+1)\cdots(\alpha_{k}+1)$.
\end{propriete}

\begin{demonstration}
Un diviseur de $n$ s'écrit $p_{1}^{\beta_{1}}\cdots p_{k}^{\beta_{k}}$ avec
$0 \leq \beta_{i} \leq \alpha_{i}$, et cette écriture est unique. Il y a donc
$\alpha_{i}+1$ choix indépendants pour chaque exposant.
\end{demonstration}

\begin{remarque}
Pour tester si un entier $n$ est premier, il suffit de chercher un diviseur
premier inférieur ou égal à $\sqrt n$ : si $n = ab$ avec $1<a\leq b$, alors
$a^{2} \leq ab = n$. Vérifier que $211$ est premier demande donc seulement
d'essayer $2$, $3$, $5$, $7$, $11$ et $13$.
\end{remarque}

\begin{entrainement}[nombres premiers]
\exo[\niveau{1}] Les entiers $91$, $101$, $143$ sont-ils premiers ?

\exo[\niveau{1}] Combien $720$ a-t-il de diviseurs positifs ?

\exo[\niveau{2}] Montrer que si $p$ est premier et $p \geq 5$, alors $p^{2}-1$
est divisible par $24$.

\exo[\niveau{3}] Montrer qu'il existe une infinité de nombres premiers.
\end{entrainement}

\begin{corrige}
\textbf{19.} $91 = 7\times13$ n'est pas premier ; $143 = 11\times13$ non plus.
En revanche $101$ n'est divisible ni par $2$, ni par $3$, ni par $5$, ni par
$7$, et $11^{2} = 121 > 101$ : il est premier.

\textbf{20.} $720 = 2^{4}\times3^{2}\times5$, donc le nombre de diviseurs
positifs vaut $(4+1)(2+1)(1+1) = 30$.

\textbf{21.} $p^{2}-1 = (p-1)(p+1)$. Comme $p$ est impair, $p-1$ et $p+1$ sont
deux entiers pairs consécutifs : l'un des deux est multiple de $4$, donc le
produit est divisible par $8$. Par ailleurs, parmi $p-1$, $p$ et $p+1$, l'un
est multiple de $3$ ; ce n'est pas $p$, qui est premier et supérieur à $3$.
Donc $3$ divise $p^{2}-1$. Comme $3$ et $8$ sont premiers entre eux, $24$
divise $p^{2}-1$.

\textbf{22.} Supposons qu'il n'y en ait qu'un nombre fini,
$p_{1},\dots,p_{r}$, et posons $N = p_{1}\times\cdots\times p_{r}+1$. L'entier
$N$ est supérieur à $1$ : il admet donc un diviseur premier, qui est l'un des
$p_{i}$. Mais ce $p_{i}$ divise aussi le produit, donc il divise leur
différence, qui vaut $1$ — ce qui est impossible. La supposition est donc
fausse.
\end{corrige}

% ===========================================================================
\section{Les théorèmes de Bézout et de Gauss}

\begin{theoreme}[théorème de Bézout]
Deux entiers relatifs $a$ et $b$ sont premiers entre eux si et seulement s'il
existe des entiers relatifs $u$ et $v$ tels que
\[
  au+bv = 1 .
\]
\end{theoreme}

\begin{demonstration}
\emph{Sens direct.} Considérons l'ensemble $E$ des entiers strictement
positifs de la forme $au+bv$. Il est non vide, donc possède un plus petit
élément $d = au_{0}+bv_{0}$. La division euclidienne $a = dq+r$ donne
\[
  r = a-dq = a\,(1-qu_{0})+b\,(-qv_{0}),
\]
qui est de la forme $au+bv$ ; comme $0 \leq r < d$ et que $d$ est le plus
petit élément strictement positif de $E$, on a $r = 0$. Donc $d \divi a$, et
de même $d \divi b$ : $d$ est un diviseur commun, donc $d = 1$ puisque $a$ et
$b$ sont premiers entre eux.

\smallskip
\emph{Réciproque.} Si $au+bv = 1$, tout diviseur commun à $a$ et $b$ divise
$1$ : le PGCD vaut $1$.
\end{demonstration}

\begin{theoreme}[théorème de Gauss]
Si $a$ divise $bc$ et si $a$ et $b$ sont premiers entre eux, alors $a$ divise
$c$.
\end{theoreme}

\begin{demonstration}
Le théorème de Bézout fournit $u$ et $v$ tels que $au+bv = 1$. En multipliant
par $c$ :
\[
  c = acu+bcv .
\]
Le premier terme est divisible par $a$ ; le second l'est aussi puisque
$a \divi bc$. Donc $a \divi c$.
\end{demonstration}

\begin{avertissement}
L'hypothèse « premiers entre eux » n'est pas décorative. $6$ divise
$4\times9 = 36$, mais $6$ ne divise ni $4$ ni $9$ : c'est que $6$ n'est
premier avec aucun des deux.
\end{avertissement}

\begin{entrainement}[Bézout et Gauss]
\exo[\niveau{1}] Montrer que $15$ et $28$ sont premiers entre eux et donner
une relation de Bézout.

\exo[\niveau{2}] Soit $n$ un entier. Montrer que $2n+1$ et $3n+2$ sont
premiers entre eux.

\exo[\niveau{2}] Montrer que si $5$ divise $3n$ alors $5$ divise $n$.

\exo[\niveau{3}] Soient $a$ et $b$ premiers entre eux. Montrer que si $a$ et
$b$ divisent $c$, alors $ab$ divise $c$.
\end{entrainement}

\begin{corrige}
\textbf{23.} $28 = 15+13$, $15 = 13+2$, $13 = 6\times2+1$ : le dernier reste
non nul est $1$, donc les deux entiers sont premiers entre eux. En remontant,
$1 = 13-6\times2 = 7\times13-6\times15 = 7\times28-13\times15$, soit
$15\times(-13)+28\times7 = 1$.

\textbf{24.} $3(2n+1)-2(3n+2) = -1$. Tout diviseur commun de $2n+1$ et de
$3n+2$ divise cette combinaison, donc divise $1$ : le PGCD vaut $1$.

\textbf{25.} $5$ divise $3n$ et $\pgcd(5\,;3) = 1$ : le théorème de Gauss
donne que $5$ divise $n$.

\textbf{26.} Comme $a$ divise $c$, écrivons $c = ak$. Alors $b$ divise $ak$
et $\pgcd(a\,;b) = 1$ : le théorème de Gauss donne que $b$ divise $k$, soit
$k = bm$. Ainsi $c = abm$, et $ab$ divise $c$.
\end{corrige}

% ===========================================================================
\section{Équations diophantiennes}

\begin{definition}[équation diophantienne du premier degré]
On appelle ainsi toute équation de la forme $ax+by = c$, où $a$, $b$, $c$ sont
des entiers donnés et où l'on cherche les couples $(x\,;y)$ d'entiers
\emph{relatifs} qui la vérifient.
\end{definition}

\begin{theoreme}[condition d'existence]
L'équation $ax+by = c$ admet des solutions entières si et seulement si
$\pgcd(a\,;b)$ divise $c$.
\end{theoreme}

\begin{demonstration}
Posons $d = \pgcd(a\,;b)$. Si $(x\,;y)$ est solution, $d$ divise $ax+by = c$.
Réciproquement, si $d \divi c$, écrivons $c = dc'$ ; le théorème de Bézout
appliqué à $\frac ad$ et $\frac bd$, qui sont premiers entre eux, donne $u$ et
$v$ tels que $\frac ad u+\frac bd v = 1$, et le couple $(uc'\,;vc')$ convient.
\end{demonstration}

\begin{methode}{résoudre $ax+by = c$}
La rédaction attendue tient en quatre temps, et le correcteur les cherche un
par un.
\begin{enumerate}[label=\textbf{\arabic*.}, leftmargin=*, itemsep=0.15em]
  \item Calculer $\pgcd(a\,;b)$ et vérifier qu'il divise $c$ — sinon,
        conclure qu'il n'y a pas de solution.
  \item Trouver une solution particulière $(x_{0}\,;y_{0})$, par
        l'algorithme d'Euclide remonté ou par essais si les nombres sont
        petits.
  \item Soustraire membre à membre l'équation générale et l'équation
        particulière, puis appliquer le théorème de Gauss.
  \item Écrire l'ensemble des solutions et vérifier qu'elles conviennent
        toutes.
\end{enumerate}
\end{methode}

\begin{exemple}[une résolution complète]
Résolvons $11x-7y = 9$ dans $\Z\times\Z$.

\smallskip
\emph{Existence.} $\pgcd(11\,;7) = 1$, qui divise $9$ : il y a des solutions.

\smallskip
\emph{Solution particulière.} On remarque $11\times2-7\times3 = 1$, donc
$11\times18-7\times27 = 9$ : le couple $(18\,;27)$ convient.

\smallskip
\emph{Solution générale.} Par différence,
$11(x-18) = 7(y-27)$. Ainsi $7$ divise $11(x-18)$ ; comme $7$ et $11$ sont
premiers entre eux, le théorème de Gauss donne $7 \divi x-18$, c'est-à-dire
$x = 18+7k$ avec $k \in \Z$. En reportant, $y = 27+11k$.

\smallskip
\emph{Vérification.} $11(18+7k)-7(27+11k) = 198+77k-189-77k = 9$. L'ensemble
des solutions est donc
\[
  \bigl\{(18+7k\,;\,27+11k) \ \mid \ k \in \Z\bigr\}.
\]
\end{exemple}

\begin{entrainement}[équations diophantiennes]
\exo[\niveau{1}] L'équation $6x+9y = 5$ a-t-elle des solutions entières ?

\exo[\niveau{2}] Résoudre dans $\Z\times\Z$ l'équation $17x-24y = 1$.

\exo[\niveau{2}] Résoudre dans $\Z\times\Z$ l'équation $5x+8y = 3$.

\exo[\niveau{3}] Un commerçant paie une somme de $2026$ ariary avec des
pièces de $50$ et de $200$ ariary. Déterminer tous les nombres possibles de
pièces de chaque sorte, ou montrer qu'il n'y en a aucun.
\end{entrainement}

\begin{corrige}
\textbf{27.} $\pgcd(6\,;9) = 3$, qui ne divise pas $5$ : l'équation n'a
aucune solution entière.

\textbf{28.} La relation de Bézout de l'exercice \textbf{17} s'écrit
$17\times(-7)-24\times(-5) = 1$ : le couple $(-7\,;-5)$ est solution
particulière. Toutes les solutions sont alors
$x = -7+24k$ et $y = -5+17k$, avec $k \in \Z$.

\textbf{29.} $5\times(-3)+8\times2 = 1$ ; en multipliant par $3$,
$5\times(-9)+8\times6 = 3$. Le couple $(-9\,;6)$ est solution particulière, et
l'ensemble des solutions est $x = -9+8k$, $y = 6-5k$, avec $k \in \Z$.

\textbf{30.} Le nombre de pièces vérifierait $50x+200y = 2026$. Le membre de
gauche est un multiple de $50$, or $2026 = 50\times40+26$ ne l'est pas : le
paiement est impossible avec ces seules pièces.
\end{corrige}

% ===========================================================================
\section{Les classes de congruence et leurs opérations}

\begin{definition}[classe de congruence]
Soit $n$ un entier naturel non nul. La \textbf{classe} de l'entier $a$ modulo
$n$, notée $\overline a$, est l'ensemble de tous les entiers congrus à $a$
modulo $n$. L'ensemble des classes est noté $\Zn n$ ; il contient exactement
$n$ éléments : $\overline 0, \overline 1, \dots, \overline{n-1}$.
\end{definition}

\begin{propriete}[opérations]
Les opérations $\overline a+\overline b = \overline{a+b}$ et
$\overline a \times \overline b = \overline{ab}$ ne dépendent pas des
représentants choisis : elles sont bien définies sur $\Zn n$.
\end{propriete}

\begin{exemple}[les tables de $\Zn 7$]
\begin{center}
\setlength{\tabcolsep}{5.5pt}
\renewcommand{\arraystretch}{1.12}
\begin{tabular}{c|ccccccc}
  $+$ & $\overline0$ & $\overline1$ & $\overline2$ & $\overline3$ & $\overline4$ & $\overline5$ & $\overline6$\\
  \hline
  $\overline0$ & $\overline0$ & $\overline1$ & $\overline2$ & $\overline3$ & $\overline4$ & $\overline5$ & $\overline6$\\
  $\overline1$ & $\overline1$ & $\overline2$ & $\overline3$ & $\overline4$ & $\overline5$ & $\overline6$ & $\overline0$\\
  $\overline2$ & $\overline2$ & $\overline3$ & $\overline4$ & $\overline5$ & $\overline6$ & $\overline0$ & $\overline1$\\
  $\overline3$ & $\overline3$ & $\overline4$ & $\overline5$ & $\overline6$ & $\overline0$ & $\overline1$ & $\overline2$\\
  $\overline4$ & $\overline4$ & $\overline5$ & $\overline6$ & $\overline0$ & $\overline1$ & $\overline2$ & $\overline3$\\
  $\overline5$ & $\overline5$ & $\overline6$ & $\overline0$ & $\overline1$ & $\overline2$ & $\overline3$ & $\overline4$\\
  $\overline6$ & $\overline6$ & $\overline0$ & $\overline1$ & $\overline2$ & $\overline3$ & $\overline4$ & $\overline5$\\
\end{tabular}
\quad
\begin{tabular}{c|ccccccc}
  $\times$ & $\overline0$ & $\overline1$ & $\overline2$ & $\overline3$ & $\overline4$ & $\overline5$ & $\overline6$\\
  \hline
  $\overline0$ & $\overline0$ & $\overline0$ & $\overline0$ & $\overline0$ & $\overline0$ & $\overline0$ & $\overline0$\\
  $\overline1$ & $\overline0$ & $\overline1$ & $\overline2$ & $\overline3$ & $\overline4$ & $\overline5$ & $\overline6$\\
  $\overline2$ & $\overline0$ & $\overline2$ & $\overline4$ & $\overline6$ & $\overline1$ & $\overline3$ & $\overline5$\\
  $\overline3$ & $\overline0$ & $\overline3$ & $\overline6$ & $\overline2$ & $\overline5$ & $\overline1$ & $\overline4$\\
  $\overline4$ & $\overline0$ & $\overline4$ & $\overline1$ & $\overline5$ & $\overline2$ & $\overline6$ & $\overline3$\\
  $\overline5$ & $\overline0$ & $\overline5$ & $\overline3$ & $\overline1$ & $\overline6$ & $\overline4$ & $\overline2$\\
  $\overline6$ & $\overline0$ & $\overline6$ & $\overline5$ & $\overline4$ & $\overline3$ & $\overline2$ & $\overline1$\\
\end{tabular}
\end{center}

Dans la table de multiplication, chaque ligne non nulle contient exactement une
fois $\overline1$ : tout élément non nul de $\Zn 7$ admet donc un inverse.
C'est une particularité des modules premiers. Dans $\Zn 6$ au contraire,
$\overline2 \times \overline3 = \overline0$ sans qu'aucun des deux facteurs
soit nul, et $\overline2$ n'a pas d'inverse.
\end{exemple}

\begin{entrainement}[classes de congruence]
\exo[\niveau{1}] Dresser les tables d'addition et de multiplication de
$\Zn 4$.

\exo[\niveau{2}] Résoudre dans $\Zn 7$ l'équation
$\overline x^{2}+\overline3\,\overline x = \overline4$.

\exo[\niveau{2}] Déterminer les éléments inversibles de $\Zn 8$.

\exo[\niveau{3}] Résoudre dans $\Zn 7$ l'équation
$\overline3\,\overline x = \overline5$, puis dans $\Zn 6$ l'équation
$\overline3\,\overline x = \overline5$. Commenter.
\end{entrainement}

\begin{corrige}
\textbf{31.} Les tables se remplissent en réduisant chaque somme et chaque
produit modulo $4$. On y lit que $\overline2 \times \overline2 = \overline0$ :
la classe $\overline2$ est un diviseur de zéro, et $\Zn4$ n'est donc pas un
corps.

\textbf{32.} L'équation s'écrit
$\overline x^{2}+\overline3\,\overline x-\overline4 = \overline0$, et le
premier membre se factorise en
$\left(\overline x-\overline1\right)\left(\overline x+\overline4\right)$.
Comme $7$ est premier, un produit nul dans $\Zn7$ impose un facteur nul :
$\overline x = \overline1$ ou $\overline x = -\overline4 = \overline3$. La
vérification donne bien $\overline1+\overline3 = \overline4$ et
$\overline2+\overline2 = \overline4$.

\textbf{33.} Une classe $\overline a$ est inversible dans $\Zn8$ si et
seulement si $\pgcd(a\,;8) = 1$ : les inversibles sont $\overline1$,
$\overline3$, $\overline5$ et $\overline7$, chacun étant son propre inverse.

\textbf{34.} Dans $\Zn7$, l'inverse de $\overline3$ est $\overline5$, car
$15 \equiv 1$. En multipliant, $\overline x = \overline5\times\overline5
= \overline{25} = \overline4$. Dans $\Zn6$ en revanche, les produits
$\overline3\,\overline x$ ne prennent que les valeurs $\overline0$ et
$\overline3$ : l'équation n'a aucune solution. La différence tient à la
primalité : dans $\Zn7$ toute classe non nulle est inversible, tandis que dans
$\Zn6$ la classe $\overline3$ est un diviseur de zéro.
\end{corrige}


% ===========================================================================
\begin{annales}
Les énoncés qui suivent sont repris de sujets réellement posés. Leur
provenance est indiquée à chaque fois. L'arithmétique se prête bien à cet
exercice : les sujets africains de série C posent, année après année, les
mêmes quatre questions que le baccalauréat Malagasy.

\exoann{Baccalauréat, Congo-Brazzaville, série C, session 2023 — exercice 1, parties 1 et 2}
\begin{enumerate}[label=\textbf{\arabic*)}, leftmargin=*, itemsep=0.15em]
  \item Démontrer, en décomposant les deux entiers en facteurs premiers, que
        $\pgcd(960\,;528) = 48$.
  \item On considère l'équation $(E) : 960u+528v = 48$, d'inconnue
        $(u\,;v) \in \Z\times\Z$.
        \begin{enumerate}[label=\textbf{\alph*)}, leftmargin=1.4em]
          \item Vérifier que $(5\,;-9)$ est solution de $(E)$.
          \item Montrer que $(E)$ équivaut à $20u+11v = 1$.
          \item Résoudre $20u+11v = 0$, puis en déduire toutes les solutions
                de $(E)$.
        \end{enumerate}
\end{enumerate}

\exoann{Baccalauréat, Congo-Brazzaville, série C, session 2023 — exercice 1, partie 3}
\begin{enumerate}[label=\textbf{\alph*)}, leftmargin=*, itemsep=0.15em]
  \item Montrer que $\congru{18x-31}{0}{7}$ équivaut à $\congru{4x}{3}{7}$.
  \item Résoudre cette congruence dans $\Z$.
\end{enumerate}

\exoann{Baccalauréat, Côte d'Ivoire, série C, session 2023 — exercice 5}
On considère l'équation $4x-y = 2$, d'inconnue $(x\,;y) \in \Z\times\Z$, et on
note $(x_{0}\,;y_{0})$ une solution.
\begin{enumerate}[label=\textbf{\arabic*)}, leftmargin=*, itemsep=0.15em]
  \item Démontrer que $y_{0}$ est pair, puis déterminer les valeurs possibles
        de $\pgcd(x_{0}\,;y_{0})$.
  \item Vérifier que $(1\,;2)$ est solution, puis résoudre l'équation.
  \item Déterminer les couples solutions dont les deux termes sont premiers
        entre eux.
\end{enumerate}

\exoann{Baccalauréat, Congo-Brazzaville, série C, session 2022 — exercice 1}
On note $(S)$ l'ensemble des couples $(x\,;y)$ d'entiers relatifs de la forme
$x = 12+3t$ et $y = 37-7t$, où $t$ décrit $\Z$.
\begin{enumerate}[label=\textbf{\arabic*)}, leftmargin=*, itemsep=0.15em]
  \item Vérifier que tout couple de $(S)$ vérifie $7x+3y = 195$.
  \item Soit $d = \pgcd(x\,;y)$. Justifier que $d$ divise $195$.
  \item Déterminer l'ensemble des diviseurs positifs de $195$.
  \item Trouver les couples de $(S)$ tels que $0 \leq y \leq 37$.
  \item En déduire le couple $(x_{0}\,;y_{0})$ pour lequel $d = 3$.
\end{enumerate}

\exoann{Baccalauréat, Bénin, série C, session 2023 — problème 1}
Soit $p$ un nombre premier strictement supérieur à $5$.
\begin{enumerate}[label=\textbf{\alph*)}, leftmargin=*, itemsep=0.15em]
  \item Démontrer que $p^{2}-1$ est divisible par $8$.
  \item Montrer que $p^{2}-1$ est un multiple de $3$.
  \item En déduire que $p^{2}-1$ est divisible par $24$.
  \item Déterminer $p^{2}$ sachant que cet entier est compris entre $3500$ et
        $4000$.
\end{enumerate}

\exoann{Baccalauréat, Bénin, série C, session 2022 — problème 3}
Déterminer les entiers naturels $a$ et $b$ tels que
\[
  \pgcd(a+b\,;ab) = 121, \qquad \ppcm(a\,;b) = 20\,757, \qquad 400 < b < a .
\]

\exoann{Baccalauréat, Madagascar, série C, session 2012 — exercice, question 1d}
Trouver un couple $(a\,;b) \in \Z\times\Z$ vérifiant
\[
  11a-7b = 9 = \pgcd(a\,;b).
\]

\exoann{Baccalauréat, Madagascar, série S, session 2022 — exercice 1, partie II}
Résoudre dans $\Zn 7$ l'équation
$\overline x^{2}+\overline3\,\overline x = \overline4$.

\exoann{Baccalauréat, Madagascar, série S, session 2022 — exercice 1, partie III}
Résoudre dans $\Z\times\Z$ l'équation $37x-43y = 1$, en utilisant l'algorithme
d'Euclide.

\exoann{Baccalauréat, Madagascar, série C, session 2026 — exercice 1}
Démontrer que, pour tout entier naturel $n$, l'entier
$N = 2^{6n+3}+3^{2n+1}$ est divisible par $11$.
\end{annales}

\begin{corrige}[annales]
\textbf{1.} $960 = 2^{6}\times3\times5$ et $528 = 2^{4}\times3\times11$ : les
facteurs communs, affectés du plus petit exposant, donnent
$2^{4}\times3 = 48$. Ensuite $960\times5+528\times(-9) = 4800-4752 = 48$. En
divisant $(E)$ par $48$, on obtient $20u+11v = 1$. L'équation homogène
$20u+11v = 0$ donne $20u = -11v$ ; comme $20$ et $11$ sont premiers entre eux,
le théorème de Gauss impose $u = 11k$ et $v = -20k$. Les solutions de $(E)$
sont donc les couples $(5+11k\,;\,-9-20k)$, $k \in \Z$.

\textbf{2.} Modulo $7$, $18 \equiv 4$ et $31 \equiv 3$, d'où l'équivalence.
Comme $4\times2 = 8 \equiv 1$, l'inverse de $\overline4$ est $\overline2$ :
$\congru{x}{2\times3}{7}$, soit $\congru{x}{6}{7}$. Les solutions sont les
entiers $x = 6+7k$.

\textbf{3.} De $y_{0} = 4x_{0}-2$ on tire que $y_{0}$ est pair. Si $d$ divise
$x_{0}$ et $y_{0}$, il divise $4x_{0}-y_{0} = 2$ : donc $d = 1$ ou $d = 2$.
Le couple $(1\,;2)$ convient. Par différence, $4(x-1) = y-2$, donc
$y = 4x-2$ : les solutions sont les couples $(k\,;4k-2)$, $k \in \Z$. Enfin
$\pgcd = 1$ exige $x$ impair, car si $x$ est pair, $2$ divise les deux termes.

\textbf{4.} $7(12+3t)+3(37-7t) = 84+21t+111-21t = 195$. Puisque $d$ divise $x$
et $y$, il divise $7x+3y = 195 = 3\times5\times13$, dont les diviseurs
positifs sont $1$, $3$, $5$, $13$, $15$, $39$, $65$ et $195$. La condition
$0 \leq 37-7t \leq 37$ donne $0 \leq t \leq 5$, soit les couples
$(12\,;37)$, $(15\,;30)$, $(18\,;23)$, $(21\,;16)$, $(24\,;9)$ et
$(27\,;2)$. Leurs PGCD valent respectivement $1$, $15$, $1$, $1$, $3$ et $1$ :
le couple cherché est $(24\,;9)$.

\textbf{5.} $p$ est impair, donc $p-1$ et $p+1$ sont deux entiers pairs
consécutifs : l'un des deux est multiple de $4$, et leur produit $p^{2}-1$ est
divisible par $8$. Par ailleurs $p$ n'est pas multiple de $3$, donc
$\congru{p}{1}{3}$ ou $\congru{p}{2}{3}$ ; dans les deux cas
$\congru{p^{2}}{1}{3}$, et $3$ divise $p^{2}-1$. Comme $8$ et $3$ sont
premiers entre eux, leur produit $24$ divise $p^{2}-1$. Enfin
$3500 \leq p^{2} \leq 4000$ donne $59 \leq p \leq 63$ ; le seul nombre premier
de cet intervalle est $61$, d'où $p^{2} = 3721$.

\textbf{6.} Tout nombre premier divisant $a+b$ et $ab$ divise $a$ et $b$ : les
facteurs premiers de $\pgcd(a+b\,;ab)$ sont ceux de $d = \pgcd(a\,;b)$. Ici
$121 = 11^{2}$, donc $11$ divise $d$ ; on pose $a = 11a'$ et $b = 11b'$ avec
$a'$ et $b'$ premiers entre eux. Alors $ab = d \times \ppcm(a\,;b)$ donne
$121\,a'b' = 11 \times 20\,757$, soit $a'b' = 1887 = 3\times17\times37$. Les
couples possibles, avec $a'+b'$ multiple de $11$, se réduisent à
$(51\,;37)$, dont la somme vaut $88$. D'où $a = 561$ et $b = 407$, qui
vérifient bien $400<407<561$.

\textbf{7.} Les solutions de $11x-7y = 9$ sont les couples
$(18+7k\,;27+11k)$. Pour $k = 0$ on obtient $(18\,;27)$, dont le PGCD vaut
$9$ : le couple $(18\,;27)$ convient.

\textbf{8.} On essaie les sept classes. Seules $\overline1$ et $\overline3$
conviennent : $1+3 = 4$ et $9+9 = 18 \equiv 4$. L'ensemble des solutions est
$\left\{\overline1\,;\overline3\right\}$.

\textbf{9.} L'algorithme d'Euclide donne $43 = 37+6$, $37 = 6\times6+1$ ; en
remontant, $1 = 37-6\times6 = 37-6(43-37) = 7\times37-6\times43$. Le couple
$(7\,;6)$ est donc solution particulière, et les solutions sont les
$(7+43k\,;\,6+37k)$, $k \in \Z$.

\textbf{10.} On écrit $2^{6n+3} = 8\times64^{n}$ et
$3^{2n+1} = 3\times9^{n}$. Or $64 = 5\times11+9$, donc
$\congru{64}{9}{11}$ et $\congru{64^{n}}{9^{n}}{11}$. Ainsi
$\congru{N}{8\times9^{n}+3\times9^{n}}{11}$, c'est-à-dire
$\congru{N}{11\times9^{n}}{11}$, donc $\congru{N}{0}{11}$.
\end{corrige}

% ===========================================================================
\begin{jecherche}[les paniers du marché d'Anosibe]
Une coopérative conditionne son riz en sacs de $23$ kg et de $17$ kg. Un
client passe commande de $1\,000$ kg exactement, et la coopérative doit lui
dire si c'est possible, et de combien de façons.

\begin{enumerate}[label=\textbf{\arabic*.}, leftmargin=*, itemsep=0.3em]
  \item Traduire le problème par une équation à deux inconnues entières.
  \item Justifier que cette équation admet des solutions dans $\Z\times\Z$.
  \item En déterminer une solution particulière.
  \item Donner l'ensemble des solutions entières.
  \item En déduire toutes les commandes possibles, sachant que les deux
        nombres de sacs doivent être des entiers naturels.
  \item La coopérative souhaite livrer le moins de sacs possible au total.
        Quelle solution doit-elle retenir ?
\end{enumerate}
\end{jecherche}

\begin{corrige}[les paniers du marché d'Anosibe]
\textbf{1.} En notant $x$ le nombre de sacs de $23$ kg et $y$ celui de sacs de
$17$ kg, la commande s'écrit $23x+17y = 1000$.

\textbf{2.} $\pgcd(23\,;17) = 1$, qui divise $1000$ : il existe des solutions
entières.

\textbf{3.} L'algorithme d'Euclide donne $23 = 17+6$, $17 = 2\times6+5$,
$6 = 5+1$ ; en remontant, $1 = 6-5 = 6-(17-2\times6) = 3\times6-17
= 3(23-17)-17 = 3\times23-4\times17$. En multipliant par $1000$ :
$23\times3000-17\times4000 = 1000$. Le couple $(3000\,;-4000)$ est solution.

\textbf{4.} Par différence, $23(x-3000) = -17(y+4000)$. Comme $23$ et $17$
sont premiers entre eux, $17$ divise $x-3000$ : $x = 3000+17k$ et
$y = -4000-23k$, avec $k \in \Z$.

\textbf{5.} Il faut $x \geq 0$ et $y \geq 0$, soit
$k \geq -176{,}47$ et $k \leq -173{,}91$ : donc $k \in \{-176,-175,-174\}$,
ce qui donne les commandes $(8\,;48)$, $(25\,;25)$ et $(42\,;2)$.

\textbf{6.} Les totaux de sacs valent respectivement $56$, $50$ et $44$. La
coopérative retiendra donc $42$ sacs de $23$ kg et $2$ sacs de $17$ kg.
\end{corrige}

% ===========================================================================
\begin{jemeteste}
Répondre par vrai ou faux, en justifiant en une ligne.

\begin{enumerate}[label=\textbf{\arabic*.}, leftmargin=*, itemsep=0.28em]
  \item Le reste de la division de $-13$ par $5$ est $-3$.
  \item Si $\congru{a}{b}{n}$ alors $\congru{a^{2}}{b^{2}}{n}$.
  \item Si $\congru{2a}{2b}{6}$ alors $\congru{a}{b}{6}$.
  \item Deux entiers consécutifs sont toujours premiers entre eux.
  \item L'équation $4x+6y = 7$ admet des solutions entières.
  \item Si un nombre premier $p$ divise $ab$, alors $p$ divise $a$ ou $p$
        divise $b$.
\end{enumerate}
\end{jemeteste}

\begin{corrige}[je me teste]
\textbf{1.} Faux — le reste est positif : $-13 = 5\times(-3)+2$, le reste vaut
$2$.
\textbf{2.} Vrai — la congruence est compatible avec le produit.
\textbf{3.} Faux — $\congru{2\times1}{2\times4}{6}$ mais
$1 \not\equiv 4 \ [6]$ ; on ne simplifie que par un entier premier avec le
module.
\textbf{4.} Vrai — tout diviseur commun à $n$ et $n+1$ divise leur différence
$1$.
\textbf{5.} Faux — $\pgcd(4\,;6) = 2$ ne divise pas $7$.
\textbf{6.} Vrai — c'est le lemme d'Euclide, conséquence du théorème de Gauss.
\end{corrige}

% ===========================================================================
\begin{commeaubac}
\bmtsource{Baccalauréat, série S, session 2022 — exercice 1, parties I à III}
\begin{enumerate}[label=\textbf{\Roman*.}, leftmargin=*, itemsep=0.3em]
  \item Démontrer que, pour tout entier naturel $n$, l'entier
        $9^{3n+1}+4^{6n+1}$ est divisible par $13$. \hfill\textup{(1 pt)}
  \item Résoudre dans $\Zn 7$ l'équation
        $\overline x^{2}+\overline3\,\overline x = \overline4$.
        \hfill\textup{(1 pt)}
  \item Résoudre dans $\Z\times\Z$ l'équation $37x-43y = 1$ en utilisant
        l'algorithme d'Euclide. \hfill\textup{(1 pt)}
\end{enumerate}
\end{commeaubac}

\begin{corrige}[comme au baccalauréat]
\textbf{I.} $9^{3} = 729 = 13\times56+1$, donc $\congru{9^{3n}}{1}{13}$ et
$\congru{9^{3n+1}}{9}{13}$. De même $4^{6} = 4096 = 13\times315+1$, donc
$\congru{4^{6n+1}}{4}{13}$. La somme est congrue à $13$, donc à $0$ : elle est
divisible par $13$.

\textbf{II.} En essayant les sept classes, on trouve $\overline1$ et
$\overline3$.

\textbf{III.} $43 = 37+6$ puis $37 = 6\times6+1$, d'où
$1 = 7\times37-6\times43$. Les solutions sont les couples
$(7+43k\,;6+37k)$, $k \in \Z$.
\end{corrige}

% ===========================================================================
\begin{concours}
\bmtsource{D'après un concours d'entrée en école d'ingénieurs}
Pour tout entier naturel $n$, on pose $u_{n} = 2^{n}+3^{n}$.

\begin{enumerate}[label=\textbf{\arabic*.}, leftmargin=*, itemsep=0.25em]
  \item Déterminer les entiers $n$ pour lesquels $u_{n}$ est divisible par
        $5$.
  \item Montrer que $\pgcd(u_{n}\,;u_{n+1})$ divise $6^{n}$.
  \item En déduire que $u_{n}$ et $u_{n+1}$ sont premiers entre eux pour tout
        $n$.
\end{enumerate}
\end{concours}

\begin{corrige}[vers le concours]
\textbf{1.} Modulo $5$ : les puissances de $2$ défilent $1,2,4,3$ et celles de
$3$ défilent $1,3,4,2$, toutes deux de période $4$. En sommant :
$u_{n} \equiv 2,\ 0,\ 3,\ 0$ selon que $n \equiv 0,1,2,3 \ [4]$. Donc $5$
divise $u_{n}$ si et seulement si $n$ est impair.

\textbf{2.} Posons $d = \pgcd(u_{n}\,;u_{n+1})$. Alors $d$ divise
$u_{n+1}-2u_{n} = 3^{n+1}-2\times3^{n} = 3^{n}$, et de même
$3u_{n}-u_{n+1} = 3\times2^{n}-2^{n+1} = 2^{n}$. Il divise donc leur produit
$6^{n}$.

\textbf{3.} Comme $d$ divise $2^{n}$ et $3^{n}$, qui sont premiers entre eux,
$d$ divise leur PGCD, égal à $1$. Donc $d = 1$ : deux termes consécutifs de la
suite sont premiers entre eux.
\end{corrige}

% ===========================================================================
\begin{notepourlaclasse}
Le chapitre paraît confortable avec ses quatorze heures ; il ne l'est pas. Les
élèves lisent les énoncés sans difficulté et se croient donc en terrain connu,
puis butent sur la rédaction. Prévoyez une séance de plus que le volume
annoncé, prise sur les révisions de cinquième période.

Deux gestes se travaillent explicitement, et rien ne les remplace : éliminer
la lettre variable dans une divisibilité, et appliquer le théorème de Gauss
après avoir soustrait la solution particulière. Faites rédiger le second
intégralement, au tableau, avant toute autonomie : c'est là que se perdent les
points, les élèves annonçant la forme générale sans jamais justifier le
passage par Gauss.

Sur les congruences, imposez le tableau des puissances successives. Un élève
qui écrit les quatre premières lignes de $2^{n}$ modulo $7$ voit la période
apparaître et n'a plus rien à mémoriser. Ceux qui refusent ce tableau
cherchent des astuces et se trompent.
\end{notepourlaclasse}

% =========================================================================
\begin{devoir}[2 heures]
Trois exercices indépendants, notés sur $20$. Le devoir est cumulatif : il porte sur ce chapitre et sur ceux qui le précèdent. Les énoncés sont repris de sessions réelles et assemblés ici en épreuve ; leur provenance est indiquée à chaque fois.

\exods{1}{Baccalauréat, Madagascar, série A, session 2026 — exercice 1}{5 points}
Soit $(U_{n})$ définie par $U_{0} = 5$ et
$U_{n+1} = \dfrac{U_{n}+3}{4}$.
\begin{enumerate}[label=\textbf{\alph*)}, leftmargin=*, itemsep=0.15em]
  \item On pose $V_{n} = U_{n}-1$. Montrer que $(V_{n})$ est géométrique et
        préciser sa raison.
  \item Exprimer $V_{n}$ puis $U_{n}$ en fonction de $n$, et donner la limite
        de $(U_{n})$.
  \item Calculer $S_{n} = U_{0}+U_{1}+\cdots+U_{n}$.
\end{enumerate}

\exods{2}{Baccalauréat, Congo-Brazzaville, série C, session 2023 — exercice 3}{7 points}
\emph{On admet que $f(x) = \frac{6\ln(x+3)}{x+3}$ est décroissante et tend
vers $0$ en $+\infty$.}

On pose $u_{n} = \displaystyle\int_{n}^{n+1}f(x)\dx$.
\begin{enumerate}[label=\textbf{\alph*)}, leftmargin=*, itemsep=0.15em]
  \item Encadrer $f(x)$ sur $\intervalle{n}{n+1}$, puis établir
        $f(n+1) \leq u_{n} \leq f(n)$.
  \item En déduire la limite de $(u_{n})$.
\end{enumerate}

\exods{3}{Baccalauréat, Congo-Brazzaville, série C, session 2023 — exercice 1, parties 1 et 2}{8 points}
\begin{enumerate}[label=\textbf{\arabic*)}, leftmargin=*, itemsep=0.15em]
  \item Démontrer, en décomposant les deux entiers en facteurs premiers, que
        $\pgcd(960\,;528) = 48$.
  \item On considère l'équation $(E) : 960u+528v = 48$, d'inconnue
        $(u\,;v) \in \Z\times\Z$.
        \begin{enumerate}[label=\textbf{\alph*)}, leftmargin=1.4em]
          \item Vérifier que $(5\,;-9)$ est solution de $(E)$.
          \item Montrer que $(E)$ équivaut à $20u+11v = 1$.
          \item Résoudre $20u+11v = 0$, puis en déduire toutes les solutions
                de $(E)$.
        \end{enumerate}
\end{enumerate}

\end{devoir}

\begin{corrige}[devoir surveillé]
\textbf{Exercice 1.} $V_{n+1} = U_{n+1}-1 = \frac{U_{n}+3}{4}-1 = \frac{U_{n}-1}{4}
= \frac14V_{n}$ : suite géométrique de raison $\frac14$, de premier terme
$V_{0} = 4$. Donc $V_{n} = 4\left(\frac14\right)^{n}$ et
$U_{n} = 1+4^{1-n}$, qui tend vers $1$. Enfin
$S_{n} = (n+1)+4\dfrac{1-\left(\frac14\right)^{n+1}}{1-\frac14}
= (n+1)+\frac{16}{3}\left(1-4^{-(n+1)}\right)$.

\textbf{Exercice 2.} Sur $\intervalle{n}{n+1}$, la décroissance donne
$f(n+1) \leq f(x) \leq f(n)$ ; l'intervalle étant de longueur $1$,
l'intégration conserve ces bornes. Comme $f(n) \to 0$, l'encadrement donne
$u_{n} \to 0$.

\textbf{Exercice 3.} $960 = 2^{6}\times3\times5$ et $528 = 2^{4}\times3\times11$ : les
facteurs communs, affectés du plus petit exposant, donnent
$2^{4}\times3 = 48$. Ensuite $960\times5+528\times(-9) = 4800-4752 = 48$. En
divisant $(E)$ par $48$, on obtient $20u+11v = 1$. L'équation homogène
$20u+11v = 0$ donne $20u = -11v$ ; comme $20$ et $11$ sont premiers entre eux,
le théorème de Gauss impose $u = 11k$ et $v = -20k$. Les solutions de $(E)$
sont donc les couples $(5+11k\,;\,-9-20k)$, $k \in \Z$.

\end{corrige}
